同策略蒸馏（OPD）~\citep{agarwal2024policy} 在大型语言模型后训练~\cite{coreteam2026mimov2flashtechnicalreport,kimiteam2026kimik3openfrontier,glm5team2026glm5vibecodingagentic} 中，将强教师模型的能力迁移至更小或能力较弱的学生模型时，正变得日益重要。与传统的离线蒸馏不同，OPD 从学生当前策略中采样轨迹，并在学生实际访问的状态上向教师进行查询。因此，所获得的监督信号可降低训练轨迹与学生自身生成行为之间的曝光偏差~\citep{lu2025onpolicydistillation}。  
目前大多数已有研究仅在单一训练领域及与训练数据高度相关的基准上评估 OPD。  
此类评估虽证实了 OPD 能提升目标任务性能，却未能区分 \emph{局部拟合} 与 \emph{更广义的策略迁移}~\citep{chu2025sft}；而这一区分对于理解 OPD 的作用机制~\citep{li2026rethinking} 及设计有效的多教师集成策略~\citep{ma2026mopd} 均至关重要。  
尤其尚不明确的是：当训练与评估在语言、推理步长或任务领域上存在差异时，OPD 的表现究竟如何。

为填补这一空白，我们开展了一项受控研究，每次仅改变一个泛化因子，同时固定其余条件，并围绕三个范围逐步扩大的问题展开。

\textbf{\emph{RQ1：OPD 对领域内分布偏移的鲁棒性如何？}}  
在固定的领域（数学）内，我们研究两个方面。首先，我们考察训练问题的相对难度~\citep{chandra2025shape,yu2026dapo,zheng2026scope}（通过教师和学生的通过率衡量）如何影响教师在领域内的性能迁移~\citep{zheng2026scope,hou2026uni}。其次，在保持领域不变但调整评估分布的情况下，分别在语言~\citep{liu2025livemathbench,sun2026olymmath}（英语$\rightarrow$ 中文）和推理跨度~\citep{lu2025r,shojaee2025the}（短跨度$\rightarrow$ 长跨度复合问题）上，考察教师性能在这些分布偏移下的迁移效果。

对于 \textbf{RQ1}，OPD 对训练问题的难度基本不敏感：教师从未解决的问题与始终能解决的问题具有同等价值，表明 OPD 传递的是教师的推理模式而非特定问题的答案。  
在评估分布偏移下，迁移能力依然有效：仅在英文短跨度数学上训练，仍能提升学生在中文和长跨度数学上的表现。

\textbf{\emph{RQ2：OPD 在不同领域间的迁移程度如何？}}  
我们探讨数学提示的监督是否能提升学生在代码和科学任务上的表现，其他领域提示能否反向迁移到数学，以及模型原点如何影响这种迁移~\citep{luo2026demystifying}。

对于 \textbf{RQ2}，OPD 能够在两个方向上将教师的能力从其训练提示的领域迁移出去，但这种迁移主要适用于 \emph{同源} 对，即能让学生在跨领域时接近教师水平；而 \emph{跨源} 对则主要在训练分布上提升学生，且迁移效果甚至不如较弱的同源教师。

\textbf{\emph{RQ3：跨领域迁移对多教师 OPD（MOPD）意味着什么？}}  
MOPD~\citep{ma2026mopd} 将每个提示路由给领域专家，看似隔离了各自贡献。我们提出疑问：考虑到每位教师的影响超出其分配领域，这种隔离是否仍然成立？

对于 \textbf{RQ3}，同样的迁移特性带来了代价：由于路由无法将教师的影响限制在其分配领域内，MOPD 中结合多位专家会导致其能力之间产生一种依赖于混合方式的 \emph{跷跷板} 效应，而非独立组合各领域的技能。  
因此，单教师 OPD 中带来优势的泛化特性，也正是使 MOPD 难以控制的原因——这是同一枚硬币的两面，而跷跷板效应为诊断 MOPD 提供了有益视角。

综合来看，我们的研究结果揭示了一个统一的图景：(1) OPD传递的是教师的推理行为，而非特定问题的解法：在领域层面，训练难度几乎无关紧要，甚至教师从未解决过的问题，其价值也与它始终能解决的问题相当。(2) 这种迁移效果受到教师与学生之间原点关系的影响。同源教师能够将学生的能力提升至接近自身水平，跨越语言、视野，甚至其他领域；而跨源教师的影响主要局限于训练数据分布，其迁移能力甚至可能不如一个较弱的同源教师。(3) 这种广泛迁移是一把双刃剑：由于每位教师都会影响其指定领域之外的能力，因此在MOPD中，提示路由无法完全隔离教师的影响。因此，专家的组合会产生一种依赖混合能力的此消彼长现象，为诊断MOPD提供了有价值的视角。